\section{Batch Accounting}

Canvass changes often enter through batch accounting rather than direct precinct
edits. Late mail ballots, provisional adjudication, duplicated damaged ballots,
and central-count batches may be accepted or rejected after election night.
RCOUNT models this with batch manifests and optional batch-specific summaries.

For every batch manifest, accepted ballots must reconcile:
\[
  \text{accepted ballots} =
  \text{counted ballots} + \text{rejected ballots}.
\]
These are declared source facts. RCOUNT verifies that the declared counts are
internally consistent; it does not decide whether a voter was eligible, whether
a cure was timely, or whether a provisional ballot should have been accepted.
For every batch summary, the summary's counted ballots must match the manifest's
counted ballots:
\[
  \text{summary counted ballots} =
  \text{manifest counted ballots}.
\]
Batch-specific summaries are adapter-dependent: a package may contain ordinary
reporting-unit summaries without batch ids, but once a summary references a
batch id, the corresponding batch manifest is required.

In the \texttt{mail-batch-added} fixture, P-001 has an election-day batch and a
late mail batch. The late mail batch is not suspicious because it is late; it is
auditable because it is declared. The event says a late-arriving mail batch was
accepted before canvass, the batch manifest declares accepted/counted/rejected
ballots, and the batch summary ties contest totals to that accounting evidence.

\begin{center}
\small
\begin{tabularx}{\textwidth}{lrrrX}
\toprule
Batch & Accepted & Counted & Rejected & Check \\
\midrule
P-001 election day & 70 & 70 & 0 & Accepted equals counted plus rejected. \\
P-001 late mail & 10 & 10 & 0 & Batch summary must count 10 ballots. \\
P-002 election day & 60 & 60 & 0 & Batch summary must count 60 ballots. \\
\bottomrule
\end{tabularx}
\end{center}

The negative \texttt{missing-batch} fixture removes the late mail manifest while
leaving a summary that references it. Arithmetic alone may still look plausible,
but the accounting evidence is missing. The verifier should therefore fail at
\texttt{batch\_summary\_total}, not at a generic package error.

This paper stops at summary and batch accounting. CVR-to-summary reconciliation,
including ballot cards, CVR rows, and parser diagnostics, belongs to V.6.
